\documentclass[10pt,a4paper]{amsart}
\usepackage[margin=19mm]{geometry}
\usepackage{amsmath,amssymb,mathtools}
\usepackage[scr=rsfs,cal=euler]{mathalfa}
\usepackage{xfrac}
\usepackage[unicode,hidelinks,bookmarks=false]{hyperref}
\newcommand{\ie}{i.e.\ }
\newcommand{\cf}{cf.\ }
\newcommand{\resp}{respectively\ }
\newcommand{\Subsection}{\subsection}
\providecommand{\nobreakdash}{\nobreak\hbox{-}}
\setlength{\emergencystretch}{2em}
\multlinegap0pt
\usepackage{amssymb}
\usepackage{xcolor}
\usepackage{mathtools}
\usepackage{algorithm}\usepackage{algpseudocode}
\newtheorem{proposition}{Proposition}
\newtheorem{remark}{Remark}
\newtheorem{theorem}{Theorem}
\newcommand{\NN}{\mathbb{N}}
\newcommand{\dr}{\mathrm{d}}
\newcommand{\dxdy}{\;\mathrm{d} x\;\mathrm{d} y}
\title{Multiple Orthogonal Polynomials on the Ball and Radial Extensions}
\author{Lidia Fernández$^\star$ \and Juan Antonio Villegas$^\star$}
\date{Source: arXiv:2606.03446v1 (CC BY 4.0)}
\begin{document}
\maketitle

\noindent\emph{Source and licence.} Multiple Orthogonal Polynomials on the Ball and Radial Extensions. Version arXiv:2606.03446v1 (CC BY 4.0), distributed by arXiv under the Creative Commons Attribution 4.0 International licence, http://creativecommons.org/licenses/by/4.0/, as recorded in the arXiv licence metadata for this exact version and shown on the abstract page https://arxiv.org/abs/2606.03446v1. Full licence notice: \texttt{licenses/arxiv-2606.03446-CC-BY-4.0.txt}.\par\smallskip

\noindent\textit{Editorial scope.} Counted unit: the article's theorem th:NNRR-2-vars (source0.tex lines 1061--1094). The article's complete original proof of it is delivered with it (source0.tex lines 1096--1152, one \texttt{proof} environment of 57 lines). No mathematics is written, repaired or re-derived: the statement and the proof are the article's own lines, in the article's own notation, and no retained line is substituted or re-flowed.  Retained uncounted as the source-local context the proof needs: lines 531--535; lines 559--566; lines 580--591; lines 662--668; lines 672--674; lines 976--979; lines 989--1004; lines 1037--1040; lines 1042--1045; lines 1049--1055.  Nothing was omitted from any retained span, so the package carries no omission record.  No retained line contains a citation, so the wrapper carries no bibliography and no bibliography entry.  No retained line contains a figure environment, an image-inclusion command or drawing code, so no image is delivered and the package declares no asset dependency.  Editorial formatting only, in addition: an AMS \texttt{amsart} preamble that restates the article's own declarations and macros used by the retained lines, namely the environments \texttt{proposition}, \texttt{remark}, \texttt{theorem} on separate counters, together with the standard packages those lines need.  The retained environments print in this document as Proposition 1, Remark 1, Theorem 1 in order of appearance, whereas the article prints the same items with the numbering of its own full document.  The complete native source of the article as distributed in the arXiv e-print for this exact version is bundled unchanged as \texttt{sources/201975/source0.tex}.
    \begin{align}
        \label{eq:norm-harmonics-2-var}
        \|Y_1^{n}(1,\theta)\|^2&= \pi(1+\delta_{n,0}) &
        \|Y_2^{n}(1,\theta)\|^2&= \pi
    \end{align}

        \begin{equation}
        \label{eq:mop-disk}
            \begin{array}{c}
            \mathbf P_{\vec j,1}^{\vec n}(\rho,\theta) =P^{(|\vec n|-2|\vec j|, \vec \mu)}_{\vec j}(\rho^2)\, \rho^{|\vec n|-2|\vec j|}\,\cos((|\vec n|-2|\vec j|)\theta), \qquad 0\leq |\vec j| \leq |\vec n|/2,
            \\[2ex]
            \mathbf P_{\vec j,2}^{\vec n}(\rho,\theta) = P^{(|\vec n|-2|\vec j|, \vec \mu)}_{\vec j}(\rho^2)\, \rho^{|\vec n|-2|\vec j|}\,\sin((|\vec n|-2|\vec j|)\theta), \qquad 0\leq |\vec j|< |\vec n|/2.
            \end{array}
        \end{equation}

    \begin{proposition}\label{prop:orthogonality-2-vars}
        %Given $\vec\mu>-1$, recall the weight functions $W_l(x,y) = (1-x^2-y^2)^{\mu_l}$ supported on the disk $D$, $l=1,2$. 
        Let $\vec n = (n_1,\dots,n_r)$ and $\vec j = (j_1,\dots,j_r)\in\NN_0^r$ be such that  $0\leq |\vec j| \leq |\vec n|/2$ and consider the polynomials $\mathbf P_{\vec j,\nu}^{\vec n}$, $\nu=1,2$, defined in \eqref{eq:mop-disk}. Then
        \begin{equation}
        \label{eq:multiple-orthogonality-disk}
            \langle\mathbf P_{\vec j,\nu}^{\vec n}(x,y),\, x^a y^b\rangle_l=\int_D \mathbf P_{\vec j,\nu}^{\vec n}(x,y)\, x^a y^b \, W_l(x,y)\dxdy =0 \qquad \text{ if }0\leq a+b< |\vec n|-2|\vec j| + 2j_l,
        \end{equation}
        for $1\leq l\leq r$ and $\nu =1,2$. Moreover, 
        $$
        \langle\mathbf P_{\vec j,\nu}^{\vec n}(x,y),\, x^a y^b\rangle_l=0 \qquad\text{ if }a+b - |\vec n|+ 2|\vec j|\text{ is odd}, \qquad 1\leq l\leq r,  \ \nu =1,2.
        $$
    \end{proposition}

        \begin{equation}
            \label{eq:Type-I-MOP-2-vars}
            \begin{aligned}
                \mathbf A_{\vec j,1,(l)}^{\vec n}(\rho,\theta)= A^{(|\vec n|-2|\vec j|, \vec \mu)}_{\vec j,l}(\rho^2)\, \rho^{|\vec n|-2|\vec j|}\, \cos((|\vec n|-2|\vec j|)\theta), \qquad 0\leq |\vec j| \leq |\vec n|/2,\\
                \mathbf A_{\vec j,2,(l)}^{\vec n}(\rho,\theta)= A^{(|\vec n|-2|\vec j|, \vec \mu)}_{\vec j,l}(\rho^2)\, \rho^{|\vec n|-2|\vec j|}\, \sin((|\vec n|-2|\vec j|)\theta), \qquad 0\leq |\vec j| < |\vec n|/2,\\
            \end{aligned}
        \end{equation}

    \begin{remark}\label{remark:degree-type-I}
    Taking into account \eqref{eq:Type-I-MOP-2-vars} and since the univariate Type~I MOP $A^{(|\vec n|-2|\vec j|, \vec \mu)}_{\vec j,l}(t)$ have degree at most $j_l-1$, the degree of the bivariate Type~I MOP $\mathbf A_{\vec j,\nu,(l)}^{\vec n}(\mathbf x)$ is at most $|\vec n|-2(|\vec j|-j_l+1)$.
    \end{remark}

    \begin{equation}
        \label{eq:type-I-function}
        \mathbf Q_{\vec j,\nu}^{\vec n}(\mathbf x)=\sum_{l=1}^r \mathbf A_{\vec j,\nu,(l)}^{\vec n}(\mathbf x)\, W_l(\mathbf x),
    \end{equation}

    \begin{proposition} \label{prop:biorthogonality}
    
    Let $\vec n = (n_1,\dots,n_r),\,\vec m=(m_1,\dots,m_r), \, \vec j = (j_1,\dots,j_r),\,\vec k =(k_1,\dots,k_r)\in\NN_0^r$ be such that $0\leq |\vec j| \leq |\vec n|/2$ and $0\leq|\vec k|\leq|\vec m|/2$. Consider the Type~II polynomial $\mathbf{P}_{\vec j,\nu}^{\vec n}$ defined in \eqref{eq:mop-disk} and the Type~I function $\mathbf Q_{\vec k,\eta}^{\vec m}$ introduced in \eqref{eq:type-I-function}. Then 
    \begin{equation}
        I:=\int_D \mathbf{P}_{\vec j,\nu}^{\vec n}(\mathbf x)\, \mathbf Q_{\vec k,\eta}^{\vec m}(\mathbf x)\, \dr \mathbf x = 0\qquad \text{ if } \nu\neq \eta \text{ or } |\vec n|-2|\vec j|\neq |\vec m|-2|\vec k|.
    \end{equation}
    Moreover, if $\nu=\eta$ and $|\vec n|-2|\vec j| = |\vec m|-2|\vec k|$, then
    \begin{equation}
        I= \begin{cases}
            0 & \text{ if }  \vec k\leq \vec j \quad \text{(componentwise)},\\
            0 & \text{ if }  |\vec j|\leq |\vec k|-2, \\
            \|Y_{\nu}^{|\vec n|-2|\vec j|}\|^2/2 & \text{ if }  |\vec j| = |\vec k| -1,
        \end{cases}
    \end{equation}
    where the squared norm of the spherical harmonic is given in \eqref{eq:norm-harmonics-2-var}.
    \end{proposition}

    \begin{equation}
        \label{eq:basis-MOP-n}
        \left\{\mathbf P_{\vec m_i,1}^{\vec n}: 0\leq i\leq |\vec n|/2\right\}\cup\left\{\mathbf P_{\vec m_i,2}^{\vec n}: 0\leq i< |\vec n|/2\right\}
    \end{equation}

    \begin{equation}
        \label{eq:basis-MOP}
        \bigcup_{k=0}^{|\vec n|}\left\{\mathbf P_{\vec m_i,1}^{\vec m_k}: 0\leq i\leq k/2\right\}\cup\left\{\mathbf P_{\vec m_i,2}^{\vec m_k}: 0\leq i< k/2\right\}
    \end{equation}

    \begin{equation}
    \label{eq:polynomial-vector-MOP}
    \begin{aligned}
        \mathbb P_{k} &:= \left(\mathbf P_{\vec m_0,1}^{\vec m_k},\dots,\mathbf P_{\vec m_{\lfloor k/2\rfloor },1}^{\vec m_k},\mathbf P_{\vec m_0,2}^{\vec m_k},\dots,\mathbf P_{\vec m_{\lfloor k/2\rfloor},2}^{\vec m_k}\right)& \text{if }k\text{ is odd,}\\
        \mathbb P_{k} &:= \left(\mathbf P_{\vec m_0,1}^{\vec m_k},\dots,\mathbf P_{\vec m_{k/2},1}^{\vec m_k},\mathbf P_{\vec m_0,2}^{\vec m_k},\dots,\mathbf P_{\vec m_{k/2-1},2}^{\vec m_k}\right)& \text{if }k\text{ is even.}
    \end{aligned}
    \end{equation}

    \begin{theorem}\label{th:NNRR-2-vars}
        Let $\vec n$ and $\vec j$ be two multi-indices such that $0\leq |\vec j|\leq|\vec n|/2$ and let $\mathbf P_{\vec j,\nu}^{\vec n}$, with $\nu\in\{1,2\}$, be associated Type~II MOPs on the disk \eqref{eq:mop-disk}. Consider $\{\vec m_k \,:\, k=0,\dots,|\vec n| \}$ a path of neighbour multi-indices such that $ |\vec m_k|=k$, $\vec m_0 = \vec 0$, $\vec m_{|\vec n|}=\vec n$, $k=0,\dots,|\vec n|$. Choose $\vec w := \vec m_{|\vec n|+1}=\vec n+\vec e_l$ for some $l\in\{1,\dots,r\}$. For each multi-index $\vec m_k$, consider the polynomial vectors $\mathbb P_k$ given in \eqref{eq:polynomial-vector-MOP}.
        
        Then, the following relations hold
        \begin{equation}
        \label{eq:NNR}
            \begin{aligned}
                x\, \mathbf P_{\vec j,\nu}^{\vec n}(\mathbf x) = \mathbf c^{(|\vec n|+1)}\mathbb P_{|\vec n|+1} + \mathbf c^{(|\vec n|)}\mathbb P_{|\vec n|} +  \sum_{k=|\vec n|-2|\vec j|-1}^{|\vec n|-1}\mathbf c^{(k)}\mathbb P_{k}, \\
                y\, \mathbf P_{\vec j,\nu}^{\vec n}(\mathbf x) = \tilde{\mathbf c}^{(|\vec n|+1)}\mathbb P_{|\vec n|+1} + \tilde{\mathbf c}^{(|\vec n|)}\mathbb P_{|\vec n|} +  \sum_{k=|\vec n|-2|\vec j|-1}^{|\vec n|-1}\tilde{\mathbf c}^{(k)}\mathbb P_{k}, 
            \end{aligned}
        \end{equation}
        assuming that the sum starts with $k=0$ if $|\vec n|-2|\vec j|=0$, and where $\mathbf c^{(k)},\,\tilde{\mathbf c}^{(k)}\in\mathbb R^{k+1}$ are given by:
        \begin{equation}
        \label{eq:NNRR-coefficients-x}
        \begin{aligned}
            \mathbf c^{(k)} &=\begin{cases}
                (c^{(k)}_{0,1},\dots,c^{(k)}_{\lfloor k/2 \rfloor ,1}, c^{(k)}_{0,1}, \dots, c^{(k)}_{\lfloor k/2 \rfloor,1}) & \text{ if }k \text{ is odd},\\
                (c^{(k)}_{0,1},\dots,c^{(k)}_{k/2,1}, c^{(k)}_{0,1}, \dots, c^{(k)}_{k/2-1,1}) & \text{ if }k \text{ is even}, 
            \end{cases} &
            c^{(k)}_{i,\nu} = \dfrac{2}{\|Y_{\nu}^{k-2i}\|^2} \langle x\, \mathbf P_{\vec j,\nu}^{\vec n}, \mathbf Q^{\vec m_{k+2}}_{\vec m_{i+1},\nu}\rangle
        \end{aligned}
        \end{equation}
        \begin{equation}
        \label{eq:NNRR-coefficients-y}
        \begin{aligned}
            \tilde{\mathbf c}^{(k)} &=\begin{cases}
                (\tilde c^{(k)}_{0,1},\dots,\tilde c^{(k)}_{\lfloor k/2 \rfloor ,1}, \tilde c^{(k)}_{0,1}, \dots, \tilde c^{(k)}_{\lfloor k/2 \rfloor,1}) & \text{ if }k \text{ is odd},\\
                (\tilde c^{(k)}_{0,1},\dots,\tilde c^{(k)}_{k/2,1}, \tilde c^{(k)}_{0,1}, \dots, \tilde c^{(k)}_{k/2-1,1}) & \text{ if }k \text{ is even}, 
            \end{cases} &
            \tilde c^{(k)}_{i,\nu} = \dfrac{2}{\|Y_{\nu}^{k-2i}\|^2} \langle y\, \mathbf P_{\vec j,\nu}^{\vec n}, \mathbf Q^{\vec m_{k+2}}_{\vec m_{i+1},\nu}\rangle.
        \end{aligned}
        \end{equation}
        
    \end{theorem}

    \begin{proof}
        In this proof, we will show the first relation, the one multiplying by $x$, as the second one is analogous.
    
        As mentioned above, each multi-index $\vec m_k$ in the path from $\vec m_0=\vec 0$ to $\vec m_{|\vec n|}$ provides us $k+1$ lineally independent polynomials of degree $k$ \eqref{eq:basis-MOP-n}, so that \eqref{eq:basis-MOP} forms a basis for the polynomials of degree $|\vec n|$.
        
        Using this argument with $\vec w=\vec n + \vec e_l$, since $|\vec w|=|\vec n|+1$ and $x \, \mathbf P_{\vec j,\nu}^{\vec n}(\mathbf x)$ is a polynomial of degree $|\vec n|+1$, we might express it as a linear combination of Type~II MOP on the disk associated to the path $\vec m_0=\vec 0,\dots,\vec m_{|\vec n|}=\vec n, \vec m_{|\vec n|+1}=\vec w$, this is
        $$
        x\, \mathbf P_{\vec j,\nu}^{\vec n}(\mathbf x) =  \sum_{k=0}^{|\vec n|+1}\mathbf c^{(k)} \mathbb P_k(\mathbf x) =  \sum_{k=0}^{|\vec n|+1}\sum_{i=0}^{k/2}\sum_{\nu=1}^2 c_{i,\nu}^{(k)} \mathbf P_{\vec m_i,\nu}^{\vec m_k}(\mathbf x),
        $$
        where we are denoting as $\mathbf c^{(k)}$ the concatenation of the vectors $(c_{i,1}^{(k)}:0\leq i\leq k/2 )$ and $(c_{i,2}^{(k)}:0\leq i< k/2 )$. 
        
        
        We will proof that $c^{(k)}_{i,\nu}=0$ whenever $k<|\vec n|-2|\vec j|-1$. Using the integral inner product with respect to the Lebesgue measure and multiplying the last expression by $\mathbf Q^{\vec m_t}_{\vec m_s,\eta}$, we have
        $$
        \langle x\, \mathbf P_{\vec j,\nu}^{\vec n}, \mathbf Q^{\vec m_t}_{\vec m_s,\eta}\rangle =  \sum_{k=0}^{|\vec n|+1}\sum_{i=0}^{k/2}\sum_{\nu=1}^2 c_{i,\nu}^{(k)} \langle  \mathbf P_{\vec m_i,\nu}^{\vec m_k},\mathbf Q^{\vec m_t}_{\vec m_s,\eta} \rangle.
        $$
        Using the biorthogonality relation from Proposition~\ref{prop:biorthogonality}, we have that $\langle  \mathbf P_{\vec m_i,\nu}^{\vec m_k},\mathbf Q^{\vec m_t}_{\vec m_s,\eta} \rangle=0$ if either $\nu\neq \eta$ or $k-2i \neq t-2 s$. Assume $\nu = \eta$ and $k-2i = t-2 s$, then $\langle \mathbf P_{\vec m_{i},\nu}^{\vec m_k}, \mathbf Q^{\vec m_t}_{\vec m_s,\nu} \rangle=0$ if $s\leq i$ (as $\vec m_s\leq\vec m_i$ iff $s\leq i$) or $|\vec m_i|=i\leq|\vec m_s|-2=s-2$. Consequently, the product vanishes unless $|\vec m_i|=i=|\vec m_s|-1=s-1$. In that case, substituting in $k-2i = t-2 s$, we get $k=t-2$. As a consequence, 
        $$
       \langle x\, \mathbf P_{\vec j,\nu}^{\vec n}, \mathbf Q^{\vec m_t}_{\vec m_s,\eta}\rangle= c^{(t-2)}_{s-1,\eta}\|Y_{\eta}^{t-2s}\|^2/2.
        $$
        From this equality, denoting $t=k+2, s=i+1$ we obtain the expressions of the coefficients in \eqref{eq:NNRR-coefficients-x}.
        
        The coefficient $c^{(t-2)}_{s-1,\eta}$ is zero whenever the product $\langle x\, \mathbf P_{\vec j,\nu}^{\vec n}, \mathbf Q^{\vec m_t}_{\vec m_s,\eta}\rangle$ is zero. Now observe
        $$
        \langle x\, \mathbf P_{\vec j,\nu}^{\vec n}, \mathbf Q^{\vec m_t}_{\vec m_s,\eta}\rangle = \langle \mathbf P_{\vec j,\nu}^{\vec n}, x\, \mathbf Q^{\vec m_{t}}_{\vec m_{s},\eta}\rangle = \sum_{l=1}^r \langle \mathbf P_{\vec j,\nu}^{\vec n}, x\, \mathbf A^{\vec m_{t}}_{\vec m_{s},\eta,(l)}\rangle_l. 
        $$
        By Type~II multiple orthogonality, see Proposition~\ref{prop:orthogonality-2-vars}, each product $\langle \mathbf P_{\vec j,\nu}^{\vec n}, x\, \mathbf A^{\vec m_{t}}_{\vec m_{s},\eta,(l)}\rangle_l$ vanishes if 
        \begin{equation}
            \label{eq:NNRR-1}
        \deg( x\, \mathbf A^{\vec m_{t}}_{\vec m_{s},\eta,(l)})<|\vec n|-2(|\vec j|-j_l).
        \end{equation}
        Recalling Remark~\ref{remark:degree-type-I},  
        \begin{equation*}
            \begin{aligned}
                \deg(x\, \mathbf A^{\vec m_{t}}_{\vec m_{s},\eta,(l)})&\leq  1 + |\vec m_t|-2(|\vec m_s|-(\vec m_s)_l+1) \\ 
                &=1+t-2(s-(\vec m_{s})_l+1)\\&=t-1-2(s-(\vec m_{s})_l),
            \end{aligned}
        \end{equation*}
        where we used $|\vec m_k|=k$.
          
        Since $(s-(\vec m_{s})_l)=|\vec m_{s}|-(\vec m_{s})_l\geq 0$, the minimum reachable value for $(s-(\vec m_{s})_l)$ is $0$, arising whenever $s=0$, in whose case $\vec m_{s}=\vec 0$. With this information, the inequality 
        \begin{equation}
            \label{eq:NNRR-12}
            \deg( x\, \mathbf A^{\vec m_{t}}_{\vec m_{s},\eta,(l)})\leq t-1-2(s-(\vec m_{s})_l) \leq t-1
        \end{equation}
        holds for every $l$.
        
        On the other hand, we have the right hand side of \eqref{eq:NNRR-1}. Since $\vec j$ might be any multi-index satisfying $0\leq |\vec j|\leq|\vec n|/2$, all we can assert is that $0\leq j_l\leq|\vec n|/2$. In this way, 
        \begin{equation}
            \label{eq:NNRR-3}
            |\vec n|-2(|\vec j|-j_l) \geq |\vec n|-2|\vec j|.
        \end{equation}
        Gathering \eqref{eq:NNRR-12} and \eqref{eq:NNRR-3}, we get that if $t-1<|\vec n|-2|\vec j|$, then \eqref{eq:NNRR-1} holds for every $l$, and consequently $c^{(t-2)}_{s-1,\eta}=0$. Equivalently, $c^{(k)}_{i,\eta}=0$ if $k<|\vec n|-2|\vec j|-1$, completing the proof.

        
        
    \end{proof}

\end{document}
